% --- 1. INFORMAZIONI GENERALI ---
\section{Informazioni Generali}
\begin{itemize}
	\item \textbf{Luogo/Piattaforma:} Server Discord
	\item \textbf{Data:} 2026-08-04
	\item \textbf{Orario di inizio:} 14:30
	\item \textbf{Orario di fine:} 15:35
	\item \textbf{Responsabile:} Filippo Idiometri
	\item \textbf{Scriba:} Sofia De Blasi
\end{itemize}

\vspace{0.5cm}
\textbf{Partecipanti presenti:}
\begin{table}[ht]
	\centering
	\renewcommand{\arraystretch}{1.2}
	\begin{tabularx}{\textwidth}{X l}
		\toprule
		\textbf{Nome e Cognome} & \textbf{Matricola} \\
		\midrule
		Sofia De Blasi & 2111014 \\
		Filippo Idiometri & 2035617 \\
		Francesco Marcon & 2110990 \\
		Armando Moda Scarati & 2082864 \\
		Anton Ricardo Rupi & 2054304 \\
		\bottomrule
	\end{tabularx}
\\ \end{table}

\vspace{0.5cm}
\textbf{Partecipanti assenti:}
\begin{table}[ht]
	\centering
	\renewcommand{\arraystretch}{1.2}
	\begin{tabularx}{\textwidth}{X l}
		\toprule
		\textbf{Nome e Cognome} & \textbf{Matricola} \\
		\midrule
		Roberto Mariano Doroftei & 2111031 \\
		\bottomrule
	\end{tabularx}
\end{table}

% --- 2. ORDINE DEL GIORNO ---
\section{Ordine del Giorno}
Gli argomenti previsti per la riunione odierna sono:
\begin{enumerate}
	\item Sprint review e retrospective.
	\item Aggiornamento dei documenti gestionali (NdP, PdP, PdQ).
	\item Presentazione e validazione dei diagrammi delle classi, package e sequenze.
	\item Pianificazione dell'avvio dello sviluppo (implementazione e test).
	\item Impostazione delle regole per il branch protection e delle checklists per le Pull Request.
	\item Assegnazione dei nuovi ruoli e pianificazione delle attività.
\end{enumerate}

% --- 3. RESOCONTO E DISCUSSIONE ---
\section{Resoconto della Riunione}

\subsection{Sprint review e retrospective}
La review ha evidenziato il completamento delle attività di studio progettuale e di aggiornamento documentale. 
In particolare, i documenti \textbf{NdP}, \textbf{PdP} e \textbf{PdQ} sono stati aggiornati in base alle decisioni 
prese nella riunione precedente. 
Durante la retrospective è emersa una difficoltà principale: la complessità nel coordinare la disponibilità di tutti 
i membri del team, che ha richiesto un maggiore impegno per allineare i lavori. Nonostante ciò, l'obiettivo dello 
sprint è stato raggiunto con successo.

\subsection{Presentazione e validazione dei diagrammi}
Il progettista ha presentato al gruppo i diagrammi delle classi, dei package e delle sequenze relativi ai flussi 
principali del sistema. Questi documenti, elaborati durante lo sprint 15, sono stati approvati all'unanimità. 
La loro realizzazione ha permesso di allineare la visione architetturale del team e costituirà la base di partenza 
per l'implementazione. I diagrammi rappresentano le firme dei metodi e le interazioni tra i componenti, fornendo 
una guida chiara per lo sviluppo.

\subsection{Pianificazione dello sprint successivo e assegnazione dei ruoli}
È stato deciso di avviare l'implementazione seguendo i diagrammi approvati. 
Per garantire la qualità del codice, lo sviluppo sarà accompagnato dalla scrittura di test unitari per le funzioni significative. 
I test dovranno rispettare le convenzioni di naming e le linee guida definite nelle Norme di Progetto.

La distribuzione del lavoro per l'implementazione è la seguente:
\begin{itemize}
	\item \textbf{Programmatori (due flussi):} Armando Moda Scarati si dedicherà allo sviluppo della logica di frontend, all'implementazione 
	delle funzionalità di esportazione e dei relativi test mentre Anton Ricardo Rupi si dedicherà all'implementazione del dominio per 
	l'integrazione LLM, allo sviluppo delle API e dei servizi di backend e relativi test.
	\item \textbf{Responsabile:} Filippo Idiometri redigerà il seguente verbale, il verbale esterno con la proponente e 
	si occuperà del Diario di Bordo.
	\item \textbf{Amministratore:} Roberto Mariano Doroftei aggiornerà \textbf{PdP} e \textbf{Pdq} ed imposterà il branch protection.
	\item \textbf{Progettista:} Sofia de Blasi inizierà la stesura della specifica tecnica con il tracciamentodel soddisfacimento dei requisiti.
	\item \textbf{Verificatore:} Francesco Marcon.
\end{itemize}

\subsection{Nuova convenzione per i commit}
È stata introdotta una nuova convenzione per i messaggi di commit per migliorare la tracciabilità:
\begin{itemize}
    \item \textbf{st}: per commit relativi alla Specifica Tecnica.
    \item \textbf{mu}: per commit relativi al Manuale Utente.
\end{itemize}
% --- 4. DECISIONI PRESE ---
\section{Decisioni Prese}

\renewcommand{\arraystretch}{1.3}
\vspace{-0.3cm}
\noindent
\begin{longtable}{c p{12cm}}
	\toprule
	\textbf{Codice} & \textbf{Decisione} \\
	\midrule
	\endfirsthead

	\toprule
	\textbf{Codice} & \textbf{Decisione} \\
	\midrule
	\endhead

	\midrule
	\multicolumn{2}{r}{\textit{Continua nella pagina successiva}} \\
	\midrule
	\endfoot

	\bottomrule
	\endlastfoot

	\textbf{VI-25.1} & {Avvio dell'implementazione seguendo le specifiche dei diagrammi, affiancata dalla scrittura di test unitari.} \\
	\textbf{VI-25.2} & {Attivazione delle branch protection sul repository principale, richiedendo PR review, test e checklist.} \\
	\textbf{VI-25.3} & {Adozione delle etichette "st" e "mu" per i commit relativi a documentazione tecnica e manuale utente.} \\
	
\end{longtable}

% --- 5. AZIONI (TODO) ---
\section{Azioni da Svolgere (To-Do)}

\renewcommand{\arraystretch}{1.3}
\begin{longtable}{p{2.3cm} c p{4.5cm} p{3cm} l}
	\toprule
	\textbf{Codice} & \textbf{Rif. Decisione} & \textbf{Descrizione Task} & \textbf{Assegnatario} & \textbf{Scadenza} \\
	\midrule
	\endfirsthead

	\toprule
	\textbf{Codice} & \textbf{Rif. Decisione} & \textbf{Descrizione Task} & \textbf{Assegnatario} & \textbf{Scadenza} \\
	\midrule
	\endhead

	\midrule
	\multicolumn{5}{r}{\textit{Continua nella pagina successiva}} \\
	\midrule
	\endfoot

	\bottomrule
	\endlastfoot

	\ghissue{341} & - & Stesura del verbale interno odierno (VI N. 25) & Filippo Idiometri & 2026-08-11 \\[4pt]
	\ghissue{342} & - & Stesura del Diario di Bordo & Armando M. Scarati & 2026-08-07 \\[4pt]
	\ghissue{338} & - & Stesura del verbale esterno (VE N. 8) & Filippo Idiometri & 2026-08-11 \\[4pt]
	\ghissue{339} & - & Aggiornamento PdP & Roberto M. Doroftei & 2026-08-11 \\[4pt]
	\ghissue{340} & - & Aggiornamento PdQ & Roberto M. Doroftei & 2026-08-11 \\[4pt]
	\ghissue{343} & VI-25.1 & Sviluppo frontend & Armando M. Scarati & 2026-08-11 \\[4pt]
	\ghissue{344} & VI-25.1 & Implementazione funzionalità di esportazione & Armando M. Scarati & 2026-08-11 \\[4pt]
	\ghissue{346} & VI-25.1 & Sviluppo api e servizi backend & Armando M. Doroftei & 2026-08-11 \\[4pt]
	\ghissue{347} & - & Stesura della specifica tecnica & Sofia de Blasi & 2026-08-11 \\[4pt]
	\ghissue{348} & VI-25.2 & Branch protection & Roberto M. Doroftei & 2026-08-11 \\[4pt]
	
\end{longtable}

% --- 6. PROSSIMA RIUNIONE ---
\section{Prossima Riunione}
\begin{itemize}
	\item \textbf{Data:} 2026-08-11
	\item \textbf{Ora:} 14:30
	\item \textbf{OdG preliminare:}
	\begin{itemize}
		\item Analisi dell'elenco delle decisioni architetturali redatto e successiva approvazione.
		\item Verifica e review del setup dell'analisi statica per la CI e degli standard di testing selezionati.
		\item Avvio della stesura dei diagrammi delle classi per la documentazione.
		\item Pianificazione del prossimo sprint e assegnazione dei ruoli.
	\end{itemize}
\end{itemize}